On considère la fonction $f$ définie sur $[0~;~+\infty[$ par~:

\[f(x) = \ln(3x^2 + 1)\]

On admet que la fonction $f$ est croissante sur $[0~;+\infty[$.

\subsection*{Partie A}

On considère la suite $(u_n)$ définie par~:

\[u_0 = 2 \quad \text{et pour tout entier } n \geqslant 0,\quad u_{n+1} = \ln \left(3u_n^2 + 1\right).\]

\begin{enumerate}
\item Calculer $u_1$.
\item La suite $(u_n)$ est représentée ci-dessous. Émettre des conjectures sur le sens de variation et la convergence de la suite $(u_n)$.

\begin{center}























\psset{xunit=0.8cm,yunit=1.6cm,arrowsize=2pt 3}
\begin{pspicture*}(-0.7,-0.4)(13,4.5)
\psgrid[gridlabels=0pt,subgriddiv=1,gridwidth=0.15pt]
\psaxes[linewidth=1.25pt,Dx=2]{->}(0,0)(0,0)(13,4.5)
\psdots(0,2)(1,2.565)(2,3.032)(3,3.353)(4,3.547)(5,3.657)(6,3.717)
(7,3.748)(8,3.765)(9,3.773)(10,3.778)(11,3.780)(12,3.781)(13,3.782)
\uput[u](12.8,0){$n$}\uput[r](0,4.4){$u_n$}
\end{pspicture*}


\end{center}

\item Démontrer par récurrence que pour tout entier naturel $n$~:

\[2 \leqslant u_n \leqslant u_{n+1} \leqslant 4.\]

\item En déduire que la suite $(u_n)$ est convergente.
\end{enumerate}

\subsection*{Partie B}

On note $\ell$ la limite de la suite $(u_n)$.

Le but de cette partie est de trouver une valeur approchée de la limite $\ell$.

On considère la fonction $g$ définie et dérivable sur $[0~;~+\infty[$ par~:

\[g(x) = \ln(3x^2+1) - x.\]

\begin{enumerate}
\item On note $g'$ la fonction dérivée de la fonction $g$.

Démontrer que pour tout réel $x \in [0~;~+\infty[$,

\[g'(x) = \dfrac{-3x^2 + 6x - 1}{3x^2 + 1}.\]

\item Démontrer que la fonction $g$ est strictement décroissante sur $[2~;~4]$.
\item Démontrer que l'équation $g(x) = 0$ admet une unique solution $\alpha$ sur $[2~;~4]$ et en donner une valeur approchée à $10^{-2}$ près.
\item Justifier que $\ell = \alpha$.
\end{enumerate}

\bigskip
